\documentclass[../../book.tex]{subfiles}

\begin{document}

Although linear algebra is out of the main scope of this course, familiarity
with its basic definitions and concepts is still essential (especially when
considering Linear PCPs (\Cref{section:r1cs}) and Lattices). In this section, we
provide a brief review of the key principles that are used throughout the
course. You can skip this section without any doubts if you are already
familiar with that, and come back to it whenever you need to refresh your
memory.

\subsection{Motivation}\label{subsec:linalg:motivation}

In this subsection, we give the answer to the question ``why one should learn linear algebra for cryptography?''
Throughout this section, fix some field $\mathbb{F}$.

\textbf{Boring answer.} In cryptography we commonly need to work with a large set of \textit{linear} relations. 
For example, the following is a basic example of a linear relation between variables $x_1,\dots,x_n$ over the field $\mathbb{F}$:
\begin{equation*}
    \alpha_1x_1 + \alpha_2x_2 + \dots + \alpha_nx_n = 0
\end{equation*} 
for some constants $\alpha_1,\dots,\alpha_n \in \mathbb{F}$. 
Since operating with individual elements in such cases becomes troublesome, we collect them into \textit{vectors} and denote them in bold: say, $\mathbf{x}=(x_1,\dots,x_n)$.
Naturally, we can perform basic operations with them such as:
\begin{equation*}
    \begin{bmatrix}
        x_1 \\ \vdots \\ x_n
    \end{bmatrix} + \begin{bmatrix}
        y_1 \\ \vdots \\ y_n
    \end{bmatrix} = \begin{bmatrix}
        x_1 + y_1 \\
        \vdots \\
        x_n + y_n
    \end{bmatrix}, \quad \lambda\begin{bmatrix}
        x_1 \\ \vdots \\ x_n
    \end{bmatrix} = \begin{bmatrix}
        \lambda x_1 \\ \vdots \\ \lambda x_n
    \end{bmatrix} \quad \textcolor{oc-gray-8}{\textit{// for some $\lambda \in \mathbb{F}$.}}
\end{equation*}

When dealing with multiple linear relations, say,
\begin{equation*}
    \alpha x_1 + \beta x_2 = 0, \quad \gamma x_1 + \delta x_3 = 0, \quad \varepsilon x_2 + \zeta x_3 = 0,
\end{equation*}
we introduce the notion of \textit{matrices}. For instance, the equation above can be more concisely written as
\begin{equation*}
    A\mathbf{x} = 0, \quad A = \begin{bmatrix}
        \alpha & \beta & 0 \\
        \gamma & 0 & \delta \\
        0 & \varepsilon & \zeta
    \end{bmatrix}, \quad \mathbf{x} = \begin{bmatrix}
        x_1 \\ x_2 \\ x_3
    \end{bmatrix}.
\end{equation*}

The ``boring'' goal of this section is to give a basic machinery: what is a matrix-vector product (\textit{i.e.}, how to compute $A\mathbf{x}$),
how to multiply two matrices, and understand other properties of matrix such as determinant or rank.

However, there is much more beauty in studying linear algebra.
For instance, surprisingly, even studying how to solve a system of linear equations has a lot of hidden geometry which is easy to visualize.
We will now try to give a more motivating and somewhat complete answer below.

\textbf{Actual answer.} The typical exposition to linear algebra indeed starts from studying \textit{systems of linear equations}. 
For instance, here is one:
\begin{equation*}
    \begin{cases}
        2x_1 + x_2 = 4 \\
        x_1 + 2x_2 = 5.
    \end{cases} \; \textcolor{oc-indigo-7}{\textit{// System (a)}}
\end{equation*}

Using straightforward manipulations, one might verify that the solution to this equation is $(x_1,x_2)=(1,2)$.
However, now consider the equation
\begin{equation*}
    \begin{cases}
        4x_1 + 2x_2 = 1 \\
        2x_1 + 1x_2 = 3.
    \end{cases} \; \textcolor{oc-indigo-7}{\textit{// System (b)}}
\end{equation*}

If you try to solve it, you will notice that this equation has no solutions. 
However, you might wonder what is fundamentally different between these equations.
Specifically, why the first equation has a solution while the second does not? 
Do we have some rule when the equation has no solutions or when it has just a single one?
Can the system of such equations have infinite solutions? (\textit{yes, it can}).
The questions can go on.

One surprising insight for analyzing these two equations is to take a look at the following values, marked in red and blue:
\begin{equation*}
    \begin{cases}
        \textcolor{oc-red-7}{2} \cdot x_1 + \textcolor{oc-blue-7}{1}\cdot x_2 = 4 \\
        \textcolor{oc-red-7}{1} \cdot x_1 + \textcolor{oc-blue-7}{2}\cdot x_2 = 5,
    \end{cases} \quad 
    \begin{cases}
        \textcolor{oc-red-7}{4}\cdot x_1 + \textcolor{oc-blue-7}{2}\cdot x_2 = 1 \\
        \textcolor{oc-red-7}{2}\cdot x_1 + \textcolor{oc-blue-7}{1}\cdot x_2 = 3.
    \end{cases}
\end{equation*}

Note that we might arrange marked quantities in vectors, yielding
\begin{equation*}
    \begin{bmatrix}
        \textcolor{oc-red-7}{2} \\
        \textcolor{oc-red-7}{1}
    \end{bmatrix}x_1 + \begin{bmatrix}
        \textcolor{oc-blue-7}{1} \\
        \textcolor{oc-blue-7}{2}
    \end{bmatrix}x_2 = \begin{bmatrix}
        4 \\ 5
    \end{bmatrix}, \quad 
    \begin{bmatrix}
        \textcolor{oc-red-7}{4} \\
        \textcolor{oc-red-7}{2}
    \end{bmatrix}x_1 + \begin{bmatrix}
        \textcolor{oc-blue-7}{2} \\
        \textcolor{oc-blue-7}{1}
    \end{bmatrix}x_2 = \begin{bmatrix}
        1 \\ 3
    \end{bmatrix}.
\end{equation*}

This way, we reduced the question from solving an abstract set of equations to: ``can we find a linear combination\footnote{Given vectors $\mathbf{v}_1,\dots,\mathbf{v}_n$, the \textit{linear combination} is an expression of form $\lambda_1\mathbf{v}_1+\dots+\lambda_n\mathbf{v}_n$.} of vectors $\begin{bmatrix}
        \textcolor{oc-red-7}{2} \\
        \textcolor{oc-red-7}{1}
    \end{bmatrix}$ and $\begin{bmatrix}
        \textcolor{oc-blue-7}{1} \\
        \textcolor{oc-blue-7}{2}
    \end{bmatrix}$ that will give $\begin{bmatrix}
        4 \\ 5
    \end{bmatrix}$?''
This problem is in fact much more easily visualizable! 
See \Cref{fig:linear-eqs-illustration}.
Notice that when two highlighted vectors are parallel, no matter which $x_1$ and $x_2$ one picks, the sum cannot be a vector that is not parallel to ones given.
Consequently, since $(1,3)$ is not parallel to $(4,2)$ and $(2,1)$, there are no solutions.
In contrast, when two vectors are not parallel, we are always guaranteed to find the solution.

\captionsetup[subfigure]{labelformat=empty}
\begin{figure}[htbp]
  \centering

  % --- LEFT SIDE: Unique Solution ---
  \begin{subfigure}[b]{0.48\textwidth}
    \centering
    
    \begin{tikzpicture}[scale=0.5, >=Stealth]
      % Grid & Axes
      \draw[gray!20, thin] (-0.5,-0.5) grid (5.5,5.5);
      \draw[->, thick, gray!70] (-0.5,0) -- (5.5,0) node[right] {$x$};
      \draw[->, thick, gray!70] (0,-0.5) -- (0,5.5) node[above] {$y$};
      
      % Ticks
      \foreach \x in {1,...,5} \draw (\x,0.05) -- (\x,-0.05) node[below, font=\tiny] {\x};
      \foreach \y in {1,...,5} \draw (0.05,\y) -- (-0.05,\y) node[left, font=\tiny] {\y};

      % Linear combination path: 1*v1 + 2*v2 = target
      \draw[dashed, oc-blue-7, thick] (2,1) -- (4,5) 
        node[midway, right] {$+2\mathbf{v}_2$};

      % Vectors
      \draw[->, ultra thick, oc-red-7] (0,0) -- (2,1) 
        node[midway, below right] {$\mathbf{v}_1$};
      \draw[->, ultra thick, oc-blue-7] (0,0) -- (1,2) 
        node[midway, above left] {$\mathbf{v}_2$};
      \draw[->, ultra thick, oc-green-5] (0,0) -- (4,5) 
        node[above left] {$\mathbf{b}$};

      % Target Point Marker
      \filldraw[oc-green-5] (4,5) circle (2.5pt);
    \end{tikzpicture}
    \caption{\textcolor{oc-indigo-7}{\textit{System (a)}}. Solution exists.}
  \end{subfigure}
  \hfill
  % --- RIGHT SIDE: Parallel / No Solution ---
  \begin{subfigure}[b]{0.48\textwidth}
    \centering

    \begin{tikzpicture}[scale=0.5, >=Stealth]
      % Grid & Axes
      \draw[gray!20, thin] (-0.5,-0.5) grid (5.5,5.5);
      \draw[->, thick, gray!70] (-0.5,0) -- (5.5,0) node[right] {$x$};
      \draw[->, thick, gray!70] (0,-0.5) -- (0,5.5) node[above] {$y$};

      % Ticks
      \foreach \x in {1,...,5} \draw (\x,0.05) -- (\x,-0.05) node[below, font=\tiny] {\x};
      \foreach \y in {1,...,5} \draw (0.05,\y) -- (-0.05,\y) node[left, font=\tiny] {\y};

      % Line spanned by parallel vectors (y = 2x)
      %\draw[dashed, gray!60, thick] (-0.25,-0.5) -- (2.75,5.5) 
      %  node[above, font=\tiny, rotate=63.4] {$\text{Span}(\mathbf{v}_1, \mathbf{v}_2)$};

      % Parallel Vectors (slightly offset side-by-side for visual clarity)
      \draw[->, ultra thick, oc-red-7] (0,0) -- (4,2) 
        node[midway, right=9pt] {$\mathbf{v}_1$};
      \draw[->, ultra thick, oc-blue-7] (0,0) -- (2,1) 
        node[midway, below right=-2pt and 0pt] {$\mathbf{v}_2$};
      
      % Target Vector
      \draw[->, ultra thick, oc-green-5] (0,0) -- (1,3) 
        node[above left] {$\mathbf{b}$};

      % Target Point Marker
      \filldraw[oc-green-5] (1,3) circle (2.5pt);
    \end{tikzpicture}
    \caption{\textcolor{oc-indigo-7}{\textit{System (b)}}. Solution does not exist.}
  \end{subfigure}

  \caption{Geometric interpretation of the system of linear equations.}
  \label{fig:linear-eqs-illustration}
\end{figure}

It turns out that similar idea is valid for system of equations in more variables.
Say, it is also easy to visualize the following three equations:
\begin{align*}
    &\textcolor{oc-indigo-7}{\textit{// System (a)}} && \textcolor{oc-indigo-7}{\textit{// System (b)}} && \textcolor{oc-indigo-7}{\textit{// System (c)}} \\
    &\begin{cases}
        2x_1 - x_3 = 1 \\
        2x_2 + x_3 = 3 \\
        x_1 + x_2 + 2x_3 = 4
    \end{cases} && \begin{cases}
        2x_1 + x_3 = 2 \\
        2x_2 + x_3 = 2 \\
        x_1 + x_2 + x_3 = 2
    \end{cases} && \begin{cases}
        3x_1 + 2x_2 + x_3 = 1 \\
        9x_1 + 6x_2 + 3x_3 = 4 \\
        6x_1 + 4x_2 + 2x_3 = 7
    \end{cases}
\end{align*}

Again, by considering these equations in the form $\mathbf{v}_1x_1 + \mathbf{v}_2x_2 + \mathbf{v}_3x_3=\mathbf{t}$ for fixed $\mathbf{v}_1,\mathbf{v}_2,\mathbf{v}_3,\mathbf{t} \in \mathbb{R}^3$, we have:
\begin{align*}
    \textcolor{oc-indigo-7}{\textit{// System (a)}}: && \begin{bmatrix}
        \textcolor{oc-red-7}{2} \\ \textcolor{oc-red-7}{0} \\ \textcolor{oc-red-7}{1}
    \end{bmatrix}x_1 + \begin{bmatrix}
        \textcolor{oc-blue-7}{0} \\ \textcolor{oc-blue-7}{2} \\ \textcolor{oc-blue-7}{1}
    \end{bmatrix}x_2 + \begin{bmatrix}
        \textcolor{oc-orange-7}{-1} \\ \textcolor{oc-orange-7}{1} \\ \textcolor{oc-orange-7}{2}
    \end{bmatrix}x_3 = \begin{bmatrix}
        1 \\ 3 \\ 4
    \end{bmatrix}, \\
    \textcolor{oc-indigo-7}{\textit{// System (b)}}: && \begin{bmatrix}
        \textcolor{oc-red-7}{2} \\ \textcolor{oc-red-7}{0} \\ \textcolor{oc-red-7}{1}
    \end{bmatrix}x_1 + \begin{bmatrix}
        \textcolor{oc-blue-7}{0} \\ \textcolor{oc-blue-7}{2} \\ \textcolor{oc-blue-7}{1}
    \end{bmatrix}x_2 + \begin{bmatrix}
        \textcolor{oc-orange-7}{1} \\ \textcolor{oc-orange-7}{1} \\ \textcolor{oc-orange-7}{1}
    \end{bmatrix}x_3 = \begin{bmatrix}
        2 \\ 2 \\ 2
    \end{bmatrix}, \\
    \textcolor{oc-indigo-7}{\textit{// System (c)}}: && \begin{bmatrix}
        \textcolor{oc-red-7}{3} \\ \textcolor{oc-red-7}{9} \\ \textcolor{oc-red-7}{6}
    \end{bmatrix}x_1 + \begin{bmatrix}
        \textcolor{oc-blue-7}{2} \\ \textcolor{oc-blue-7}{6} \\ \textcolor{oc-blue-7}{4}
    \end{bmatrix}x_2 + \begin{bmatrix}
        \textcolor{oc-orange-7}{1} \\ \textcolor{oc-orange-7}{3} \\ \textcolor{oc-orange-7}{2}
    \end{bmatrix}x_3 = \begin{bmatrix}
        1 \\ 4 \\ 7
    \end{bmatrix}.
\end{align*}

As before, we can sketch how vectors in \textcolor{oc-red-7}{red}, \textcolor{oc-blue-7}{blue}, and \textcolor{oc-orange-7}{orange} are positioned with respect to each other.
Specifically, consider \Cref{fig:linear-eqs-3d-illustration}.
Notice that there are three primary cases how such vectors and righthand side are located in the space:
\begin{itemize}
    \item Vectors lie on the same \textit{line} (as illustrated by \textit{System (c)}).
    Since $(1,4,7)$ does not lie on the same line as the other vectors, the system has no solutions.
    Otherwise, it has an infinite number of solutions.
    \item Vectors lie on the same \textit{plane} (as illustrated by \textit{System (b)}).
    Since $(2,2,2)$ lies on this plane, there is an infinite number of solutions.
    Specifically, this set is given by $\{(1-\frac{1}{2}t,1-\frac{1}{2}t,t): t \in \mathbb{R}\}$.
    \item Vectors are lying ``arbitrarily'' (neither on the same line nor plane).
    In such case, the system of equations always has a single \textit{unique} solution.
\end{itemize}

Notice that the three-dimensional case is significantly more complicated than the two-dimensional.
Moreover, notice that both \textit{System (b)} and \textit{System (c)} might have an infinite number of solutions, yet these numbers are in a way ``different''.
Specifically, as we observed, we only need a single ``free'' variable $t \in \mathbb{R}$ to describe the set of solutions of \textit{System (b)}.
In turn, if \textit{System (c)} had solutions, we would need two free variables to describe it. 

\begin{figure}[htbp]
  \centering

  % --- ROW (a): Unique Solution ---
  \begin{subfigure}[c]{0.45\textwidth}
    \centering
    \includegraphics[width=\linewidth]{contents/images/linear-algebra/3d-system-1.png}
  \end{subfigure}%
  \hfill
  \begin{subfigure}[c]{0.50\textwidth}
    \caption{\textcolor{oc-indigo-7}{\textit{System (a)}} \\[2pt]
             Three vectors situated arbitrarily in space (as we will further formally define, they \textit{span the whole $\mathbb{R}^3$}).
             The solution to the equation always exists regardless of the righthand side.}
  \end{subfigure}

  \vspace{1em}

  % --- ROW (b): Infinite Solutions (1 Free Var / Coplanar) ---
  \begin{subfigure}[c]{0.45\textwidth}
    \centering
    \includegraphics[width=\linewidth]{contents/images/linear-algebra/3d-system-2.png}
  \end{subfigure}%
  \hfill
  \begin{subfigure}[c]{0.50\textwidth}
    \caption{\textcolor{oc-indigo-7}{\textit{System (b)}} \\[2pt]
             Three vectors are situated on the same plane.
             Consequently, the equation will have an infinite number of solutions if the righthand side also lies on this plane.
             Otherwise, the system has no solutions.}
  \end{subfigure}

  \vspace{1em}

  % --- ROW (c): Infinite Solutions (2 Free Vars / Collinear) ---
  \begin{subfigure}[c]{0.45\textwidth}
    \centering
    \includegraphics[width=\linewidth]{contents/images/linear-algebra/3d-system-3.png}
  \end{subfigure}%
  \hfill
  \begin{subfigure}[c]{0.50\textwidth}
    \caption{\textcolor{oc-indigo-7}{\textit{System (c)}}\\[2pt]
             Three vectors situated on the same line. Solution only exists if the righthand side lies along this line.}
  \end{subfigure}

  \caption{Geometric interpretation of $3 \times 3$ linear systems over $\mathbb{R}$ given by $\mathbf{v}_1x_1 + \mathbf{v}_2x_2 + \mathbf{v}_3x_3 = \mathbf{t}$ for the righthand side $\mathbf{t} \in \mathbb{R}^3$.}
  \label{fig:linear-eqs-3d-illustration}
\end{figure}

Now, here is the thing: in cryptography, we typically need to operate over systems of \textit{thousands} of linear relations.
While we could ``manually'' analyze simple cases as $2 \times 2$ or $3 \times 3$ linear equations, we need a more systematic way to give answers to the preceding questions.

The linear algebra gives all the necessary tools. 
Specifically, as you could have already noticed, the key of studying such equations is to see how column-vectors of some equation are related.
It is then convenient to collect all such vectors in a single object called a \textit{matrix}.
For instance, for \textit{System (a)}, we have:
\begin{equation*}
    A\mathbf{x} = \mathbf{t}, \quad A = \begin{bmatrix}
        2 & 0 & -1 \\
        0 & 2 & 1 \\
        1 & 1 & 2
    \end{bmatrix}, \quad \mathbf{t} = \begin{bmatrix}
        1 \\ 3 \\ 4
    \end{bmatrix}.
\end{equation*}

While of course the number of solutions depends on the righthand side, the primary structure of the solution set is \emph{determined by a map $\mathbf{x} \mapsto A\mathbf{x}$}.
In a sense, this is exactly what we are going to study in this section.
For example, when we develop the whole toolkit, here is how we can answer the previous questions systematically:

\begin{proposition}\label{prop:basic-facts}
    Consider the linear equation $A\mathbf{x}=\mathbf{t}$ for arbitrary matrix $A \in \mathbb{R}^{n \times n}$ and $\mathbf{t} \in \mathbb{R}^n$. Then,
    \begin{enumerate}
    \item there is always a single solution to the equation iff $\det A \neq 0$;
    \item \textit{(Rouché-Capelli Theorem)} the solution exists iff $\text{rank}(A \| \mathbf{t})=\text{rank}(A)$;
    \item if the solution exists, the number of variables needed to describe the solution set is $n-\text{rank}(A)$.
    \end{enumerate}
\end{proposition}

Besides, the tools we provide are as geometric as the motivation above. 
For example, $\det A$ represents the volume of the parallelepiped spanned by columns of $A$ (see figure below).
So the statement \textit{(i)} of \Cref{prop:basic-facts} basically tells: ``if columns of $A$ span non-zero volume, we can always find a unique solution.''
This fully coincides with the situation we observed in \Cref{fig:linear-eqs-3d-illustration}: line and plane technically have no volume (determinant is $0$), so the equation has either no or infinite number of solutions.
\begin{figure}[ht]
\centering
\begin{tikzpicture}[scale=0.75, >=Stealth]

% ============================================================
% Common style:
%   axes: thick gray
%   vectors: ultra thick colored arrows
% ============================================================


% ============================================================
% 1. det(A) != 0 : three-dimensional parallelepiped
% ============================================================
\begin{axis}[
    at={(0,0)},
    width=5.2cm,
    height=5.2cm,
    view={120}{25},
    axis lines=center,
    axis line style={thick, gray!70},
    xmin=0, xmax=4.2,
    ymin=0, ymax=3.8,
    zmin=0, zmax=3.5,
    ticks=none,
    xlabel style={
        at={(axis description cs:1.02,0.03)},
        anchor=west
    },
    ylabel style={
        at={(axis description cs:0.02,1.03)},
        anchor=south
    },
    zlabel style={
        at={(axis description cs:0.05,0.98)},
        anchor=south
    },
]

% Parallelepiped
\addplot3[
    fill=blue!20,
    draw=blue!60,
    opacity=0.75
] coordinates {
    (0,0,0) (3,0,0) (3,2,0) (0,2,0) (0,0,0)
};

\addplot3[
    fill=blue!20,
    draw=blue!60,
    opacity=0.75
] coordinates {
    (0,0,2) (3,0,2) (3,2,2) (0,2,2) (0,0,2)
};

\addplot3[blue!60, opacity=0.75] coordinates {(0,0,0) (0,0,2)};
\addplot3[blue!60, opacity=0.75] coordinates {(3,0,0) (3,0,2)};
\addplot3[blue!60, opacity=0.75] coordinates {(3,2,0) (3,2,2)};
\addplot3[blue!60, opacity=0.75] coordinates {(0,2,0) (0,2,2)};

% Vectors
\addplot3[
    ->,
    ultra thick,
    oc-red-7
] coordinates {(0,0,0) (3,0,0)};

\addplot3[
    ->,
    ultra thick,
    oc-blue-7
] coordinates {(0,0,0) (0,2,0)};

\addplot3[
    ->,
    ultra thick,
    oc-green-7
] coordinates {(0,0,0) (0,0,2)};

\node[
    oc-red-7,
    font=\small
] at (axis cs:1.5,-0.35,0.2) {$\mathbf{v}_1$};

\node[
    oc-blue-7,
    font=\small
] at (axis cs:-0.7,1,0.2) {$\mathbf{v}_2$};

\node[
    oc-green-7,
    font=\small
] at (axis cs:-0.25,0.3,1.15) {$\mathbf{v}_3$};

\node[
    font=\large\bfseries,
    anchor=south
] at (axis description cs:0.5,1.08)
{$\det(A)\neq0$};

\node[
    font=\small,
    anchor=north
] at (axis description cs:0.5,-0.15)
{3-dimensional span};

\end{axis}


% ============================================================
% 2. Plane : det(A) = 0
% ============================================================
\begin{axis}[
    at={(6cm,0)},
    width=5.2cm,
    height=5.2cm,
    view={120}{25},
    axis lines=center,
    axis line style={thick, gray!70},
    xmin=0, xmax=4.2,
    ymin=0, ymax=3.8,
    zmin=0, zmax=3.5,
    ticks=none,
    xlabel style={
        at={(axis description cs:1.02,0.03)},
        anchor=west
    },
    ylabel style={
        at={(axis description cs:0.02,1.03)},
        anchor=south
    },
    zlabel style={
        at={(axis description cs:0.05,0.98)},
        anchor=south
    },
]

\addplot3[
    surf,
    domain=0:5,
    y domain=0:5,
    samples=8,
    samples y=8,
    shader=flat,
    draw=gray!55,
    fill=gray!15,
    opacity=0.45,
    line join=round
]
({x},{y},{0});
% Three coplanar vectors
\addplot3[
    ->,
    ultra thick,
    oc-red-7
] coordinates {(0,0,0) (2,0,0)};

\addplot3[
    ->,
    ultra thick,
    oc-blue-7
] coordinates {(0,0,0) (0,2,0)};

\addplot3[
    ->,
    ultra thick,
    oc-green-7
] coordinates {(0,0,0) (3,2,0)};

\node[
    oc-red-7,
    font=\small
] at (axis cs:1.5,-0.15,0.5) {$\mathbf{v}_1$};

\node[
    oc-blue-7,
    font=\small
] at (axis cs:-0.3,2,0.35) {$\mathbf{v}_2$};

\node[
    oc-green-7,
    font=\small
] at (axis cs:2,2,0) {$\mathbf{v}_3$};

\node[
    font=\large\bfseries,
    anchor=south
] at (axis description cs:0.5,1.08)
{$\det(A)=0$};

\node[
    font=\small,
    anchor=north
] at (axis description cs:0.5,-0.15)
{2-dimensional span (plane)};

\end{axis}


% ============================================================
% 3. Line : det(A) = 0
% ============================================================
\begin{axis}[
    at={(12cm,0)},
    width=5.2cm,
    height=5.2cm,
    view={120}{25},
    axis lines=center,
    axis line style={thick, gray!70},
    xmin=0, xmax=5.2,
    ymin=0, ymax=3.8,
    zmin=0, zmax=4.2,
    ticks=none,
    xlabel style={
        at={(axis description cs:1.02,0.03)},
        anchor=west
    },
    ylabel style={
        at={(axis description cs:0.02,1.03)},
        anchor=south
    },
    zlabel style={
        at={(axis description cs:0.05,0.98)},
        anchor=south
    },
]

\addplot3[
    very thick,
    dashed,
    gray
] coordinates {
    (0,0,0)
    (3,3,0)
};

\addplot3[
    very thick,
    dashed,
    gray
] coordinates {
    (3,3,3)
    (3,3,0)
};

% Three collinear vectors -- made longer
\addplot3[
    ->,
    ultra thick,
    oc-green-7
] coordinates {(0,0,0) (4,4,4)};

\addplot3[
    ->,
    ultra thick,
    oc-blue-7
] coordinates {(0,0,0) (2.75,2.75,2.75)};

\addplot3[
    ->,
    ultra thick,
    oc-red-7
] coordinates {(0,0,0) (1.5,1.5,1.5)};

\node[
    oc-green-7,
    font=\small
] at (axis cs:4,3.75,4.4) {$\mathbf{v}_3$};

\node[
    oc-blue-7,
    font=\small
] at (axis cs:2.75,2.5,3.25) {$\mathbf{v}_2$};

\node[
    oc-red-7,
    font=\small
] at (axis cs:1.5,1.25,2) {$\mathbf{v}_1$};

\end{axis}

\end{tikzpicture}
\end{figure}

\begin{remark}
    In areas such as Machine Learning, studying large linear systems might give much insight into the structure of the underlying data.
    Toolkit needed for such areas is arguably much more advanced than what we will require in our course: for instance, it requires analyzing the \textit{spectrum} of linear operators.
    That said, this section focuses on the very basics.
\end{remark}

\subsection{Vector Space and Linear Subspace}

Similarly to group theory working with \textit{groups}, the linear algebra also
has a special designated primitive --- \textbf{vector space}. If previously we
were working with the (finite) field $\mathbb{F}$, now we will work with the
vector space $V$ over this field. In many \textit{practical} applications, vector space
is formed by \textbf{vectors} consisting of a finite fixed collection of
elements from the field $\mathbb{F}$. 
The most common example we will use is $\mathbb{F}^n$: the set of all $n$-tuples $(x_1,x_2,\dots,x_n)$ of
elements from $\mathbb{F}$.
We will use bold letters $\mathbf{x}$ to denote such a collection for brevity.

Let us now give a bit more general definition.

\begin{definition}
    A \textbf{vector space} $V$ over the field $\mathbb{F}$ is an abelian group for addition $+$ together with a scalar multiplication operation $\cdot$ from $\mathbb{F} \times V$ to $V$, sending $(\lambda,\mathbf{x}) \mapsto \lambda \mathbf{x}$ and such that for any $\mathbf{v},\mathbf{u} \in V$ and $\lambda,\mu \in \mathbb{F}$ we have:
    \begin{xmulticols}{2}
        \begin{itemize}
            \item $\lambda(\mathbf{u}+\mathbf{v}) = \lambda \mathbf{u} + \lambda \mathbf{v}$
            \item $(\lambda + \mu)\mathbf{v} = \lambda \mathbf{v} + \mu \mathbf{v}$
            \item $(\lambda \mu)\mathbf{v} = \lambda(\mu \mathbf{v})$
            \item $1\cdot \mathbf{v} = \mathbf{v}$
        \end{itemize}
    \end{xmulticols}

    Any element $\mathbf{v} \in V$ is called a \textbf{vector}, and any element $\lambda \in \mathbb{F}$ is called a \textbf{scalar}. We also mark vector elements in boldface.
\end{definition}

\vspace{-4mm}

\begin{example}
    For example, $V = \mathbb{F}^n$ with operations defined as:
    \begin{xequation*}
        \begin{aligned}
            \lambda \cdot (x_1,x_2,\dots,x_n) &= (\lambda x_1, \lambda x_2, \dots, \lambda x_n) \\
            (x_1,x_2,\dots,x_n) + (y_1,y_2,\dots,y_n) &= (x_1+y_1, x_2+y_2, \dots, x_n+y_n)
        \end{aligned}
    \end{xequation*}
    is a vector space. Similarly, the following three sets $V_1,V_2,V_3$ with operations defined above are also valid vector spaces:
    \begin{xequation*}
        \begin{aligned}
            &V_1 := \{(x_1,x_2,\dots,x_n) \in \mathbb{F}^n: x_1 = 0\} \\
            &V_2 := \{(x_1,x_2,\dots,x_n) \in \mathbb{F}^n: x_1-x_3 = 0\} \\
            &V_3 := \{(x_1,x_2,\dots,x_n) \in \mathbb{F}^n: x_1+x_2+\dots+x_n = 0\}
        \end{aligned}
    \end{xequation*}

    \textbf{Reasoning.} Let us see why, for example, $V_3$ is a valid vector space. Since all operations 
    are inherited from $\mathbb{F}^n$, we only need to check that $V_3$ is closed under addition and scalar multiplication.
    Suppose $\mathbf{u}=(u_1,\dots,u_n) \in V_3$ and $\mathbf{w}=(w_1,\dots,w_n) \in V_3$. Then,
    $\mathbf{u}+\mathbf{w}=(u_1+w_1,\dots,u_n+w_n)$.
    Note that $\sum_{i=1}^n u_i+w_i = \left(\sum_{i=1}^n u_i\right) + \left(\sum_{i=1}^nw_i\right) = 0 + 0 = 0$, so indeed $\mathbf{u}+\mathbf{w} \in V_3$.
    Similarly one can check that $\lambda\mathbf{u} \in V_3$ for any $\lambda \in \mathbb{F}$.
\end{example}

\begin{remark}
    Although practically we will only encounter vector spaces of the form $\mathbb{F}^n$, the notion of vector space is very broad. For example, consider the following two sets:
    \begin{align*}
        U_1 &:= \mathbb{R}[X]^{\leq n}, \; \textcolor{gray}{\textit{// polynomials in $\mathbb{R}$ of maximum degree $n$}} \\
        U_2 &:= \{\text{all functions $f: [0,1] \to \mathbb{R}$ continuous over $[0,1]$}\}.
    \end{align*}
    One can check, that both $U_1$ and $U_2$ are in fact vector spaces over $\mathbb{R}$ with $U_1$ being of finite dimension while $U_2$ being infinitely-dimensional.
\end{remark}

Very common in practice we need to work with vector spaces that are subspaces of $\mathbb{F}^n$.
Naturally, such spaces have a dedicated term.

\begin{definition}
    Let $V$ be a vector space over $\mathbb{F}$ and let $W \subseteq V$.
    $W$ is called a \textbf{linear subspace} of $V$ if $W$ is itself a vector space with operations inherited from $V$.
    Equivalently, for any $\mathbf{w}_1,\mathbf{w}_2 \in W$ and $\alpha,\beta \in \mathbb{F}$ one has
    \begin{equation*}
        \alpha\mathbf{w}_1 + \beta\mathbf{w}_2 \in W.
    \end{equation*}
\end{definition}

\begin{example}
    Let $n \geq 2$ and let
    \begin{xequation*}
        \begin{aligned}
            U &:= \{(x_1,x_2,\dots,x_n) \in \mathbb{F}^n: x_1 = x_2 = 0\} \\
            W &:= \{(x_1,x_2,\dots,x_n) \in \mathbb{F}^n: x_1 + x_2 = 0\}.
        \end{aligned}
    \end{xequation*}
    Then, \textit{(i)} both $U$ and $W$ are linear subspaces of $\mathbb{F}^n$ and \textit{(ii)} $U$ is a vector subspace of $W$.
\end{example}

\begin{example}\label{ex:span}
    One essential type of vector spaces is the so-called \textit{span}. 
    Let $V$ be a vector space over $\mathbb{F}$ and let $\mathbf{v}_1,\dots,\mathbf{v}_n \in V$. Define
    \begin{equation*}
        \text{span}(\mathbf{v}_1,\dots,\mathbf{v}_n) := \{\lambda_1\mathbf{v}_1 + \dots + \lambda_n\mathbf{v}_n: \lambda_1,\dots,\lambda_n \in \mathbb{F}\}.
    \end{equation*}
    Then, $\text{span}(\mathbf{v}_1,\dots,\mathbf{v}_n)$ is a linear subspace of $V$.
\end{example}

One important example of a vector subspace that we considered in the motivation part of this section is \textit{the set of solutions of a homogeneous\footnote{``homogeneous'' is a fancy way to say that the righthand side of an equation is $0$.} system of linear equations}.
\begin{exercise}
    Let $V$ be the set of vectors $(x_1,\dots,x_n)$ in $\mathbb{F}^n$ that satisfy:
    \begin{equation*}
        \mathbf{u}_1x_1 + \dots + \mathbf{u}_nx_n = 0, \; \text{for some $\mathbf{u}_1,\dots,\mathbf{u}_n \in \mathbb{F}^m$}
    \end{equation*}
    Show that $V$ is the linear subspace of $\mathbb{F}^n$.

    \textbf{Hint}: if $\mathbf{x}, \mathbf{y} \in \mathbb{F}^n$ are two solutions of such equation (\textit{i.e.}, $\mathbf{x},\mathbf{y} \in V$), what can we say about $\lambda\mathbf{x}$ and $\mathbf{x}+\mathbf{y}$?
\end{exercise}

\subsection{Basis and Dimensionality}

Now we define one very important notion: the linear dependence and independence of vectors.
This will allow us to define the \textit{dimensionality} of vector spaces which will describe their capacity or degree of freedom.

\begin{definition}
    A set of vectors $\mathbf{v}_1, \mathbf{v}_2, \dots, \mathbf{v}_n$ in a vector space $V$ over a field $\mathbb{F}$ is said to be \textbf{linearly independent} if the only scalars $\lambda_1, \lambda_2, \dots, \lambda_n \in \mathbb{F}$ that satisfy the equation
    \begin{equation*}
        \lambda_1 \mathbf{v}_1 + \lambda_2 \mathbf{v}_2 + \cdots + \lambda_n \mathbf{v}_n = \mathbf{0}
    \end{equation*}
    are $\lambda_1 = \lambda_2 = \cdots = \lambda_n = 0$. If there exist scalars, not all zero, that satisfy this equation, then the vectors are said to be \textbf{linearly dependent}.
\end{definition}

\begin{example}
    Consider the vectors $\mathbf{v}_1 = (1, 0, 0)$, $\mathbf{v}_2 = (0, 1, 0)$, and $\mathbf{v}_3 = (0, 0, 1)$ in $\mathbb{R}^3$. These vectors are linearly independent because the only solution to
    \begin{equation*}
        \lambda_1 \mathbf{v}_1 + \lambda_2 \mathbf{v}_2 + \lambda_3 \mathbf{v}_3 = \mathbf{0}
    \end{equation*}
    is $\lambda_1 = \lambda_2 = \lambda_3 = 0$. 
\end{example}

\begin{example}
    Consider the vectors $\mathbf{u}_1 = (1, 2, 3)$, $\mathbf{u}_2 = (2, 4, 6)$ in $\mathbb{R}^3$. These vectors are linearly dependent because $\mathbf{u}_2 = 2\mathbf{u}_1$, so the equation
    \begin{equation*}
        \lambda_1 \mathbf{u}_1 + \lambda_2 \mathbf{u}_2 = \mathbf{0}
    \end{equation*}
    has non-trivial solutions, such as $\lambda_1 = 2$ and $\lambda_2 = -1$.
\end{example}

We now describe the notion of \textit{basis}. 
Roughly, this is a set of vectors in $V$ such that \textit{(i)} they are linearly independent and \textit{(ii)} every vector in $V$ can be expressed as a linear combination of them.
More formally, we have the following definition.
\begin{definition}
    Vectors $\mathbf{e}_1,\dots,\mathbf{e}_n$ are a \textbf{basis} of $V$ if:
    \begin{itemize}
        \item $\text{span}(\mathbf{e}_1,\dots,\mathbf{e}_n) = V$ (see \Cref{ex:span} for the definition of $\text{span}(\cdot)$);
        \item $\mathbf{e}_1,\dots,\mathbf{e}_n$ are linearly independent.
    \end{itemize}
\end{definition}

This definition requires some examples.

\begin{example}
    For the vector space $\mathbb{F}^2$, vectors $\mathbf{e}_1=(1,0)$ and $\mathbf{e}_2=(0,1)$ form a basis.
    Indeed, they are linearly independent and any vector $(x_1,x_2) \in \mathbb{F}^2$ can be written as:
    \begin{equation*}
        \begin{bmatrix}
            x_1 \\ x_2
        \end{bmatrix} = x_1\begin{bmatrix}
            1 \\ 0
        \end{bmatrix} + x_2\begin{bmatrix}
            0 \\ 1
        \end{bmatrix} = x_1\mathbf{e}_1 + x_2\mathbf{e}_2.
    \end{equation*}
    Such basis is commonly referred to as \textit{canonical}. 
\end{example}

\begin{example}
    Basis is typically not unique. 
    For instance, for $\mathbb{R}^2$, vectors $\mathbf{u}_1=(1,2)$ and $\mathbf{u}_2=(2,1)$ also form a basis.
    Indeed, one can first check that they are linearly independent. 
    Besides, for any $(x_1,x_2) \in \mathbb{R}$, one can check:
    \begin{equation*}
        \begin{bmatrix}
            x_1 \\ x_2
        \end{bmatrix} = \left(-\frac{1}{3}x_1 + \frac{2}{3}x_2\right)\mathbf{u}_1 + \left(\frac{2}{3}x_1 - \frac{1}{3}x_2\right)\mathbf{u}_2,
    \end{equation*}
    so that any vector of $\mathbb{R}^2$ is expressed as some linear combination of $\mathbf{u}_1$ and $\mathbf{u}_2$.
    Moreover, any two linearly independent vectors in $\mathbb{F}^2$ are a valid basis of $\mathbb{F}^2$!
\end{example}

\begin{example}
    Recall that $\mathbb{F}[x]^{\leq n}$ (the space of polynomials with a degree at most $n$) is a vector space over $\mathbb{F}$.
    Its canonical basis can be described as:
    \begin{equation*}
        e_1 = 1, \; e_2 = x, \; e_3 = x^2, \; \dots \; e_{n+1} = x^n.
    \end{equation*}
    Indeed, $e_1,\dots,e_{n+1}$ are linearly independent (otherwise the equation $c_0+c_1x+\dots+c_nx^n=0$ could hold identically for some non-trivial choice of constants $c_0,\dots,c_n \in \mathbb{F}$).
    Clearly these elements span the whole $\mathbb{F}[x]^{\leq n}$.
\end{example}

\begin{example}
    Again, we stress on the fact that the basis might not be unique.
    For instance, for $\mathbb{F}_{5}[x]^{\leq 2}$, the basis could be 
    \begin{equation*}
        e_1 = 3, \; e_2 = x + 2, \; e_3 = x^2 + 4.
    \end{equation*}
\end{example}

Notice that while the basis in not unique for the given vector space, the number of basis vectors was the same.
This raises the natural question: can we find the basis $\mathbf{e}_1,\dots,\mathbf{e}_m$ for $\mathbb{F}^n$ or $\mathbb{F}[x]^{<n}$ with $m \neq n$?
The answer is no.

\begin{lemma}
    Any two bases of a vector space have the same length.
\end{lemma}

This motivates us to give the following definition.

\begin{definition}
    The \textbf{dimension} of the vector space $V$, denoted by $\dim V$, is the length of any basis of $V$.
\end{definition}

\begin{example}
    As we have observed earlier, $\dim \mathbb{F}^n=\dim \mathbb{F}[x]^{<n} = n$.
\end{example}

\subsection{Matrix}

Besides vectors, frequently we are working with \textbf{matrices}. In the most
basic sense, the matrix is a rectangular array of numbers, symbols, or
expressions, arranged in rows and columns. For example, the matrix $A$ with $m$
rows and $n$ columns, consisting of elements from the field $\mathbb{F}$
is denoted as $A \in \mathbb{F}^{m \times n}$. Additionally, we use notation $A
= (a_{i,j})_{i,j=1}^{m \times n}$ to denote the square matrix $A$ of size $m
\times n$ with elements $a_{i,j}$. Now, let us define operations on matrices.

\begin{definition}[Matrix operations]\label{def:matrix-ops}
    We define the following operations over matrices:
    \begin{itemize}
        \item for matrices $A,B \in \mathbb{F}^{m \times n}$ of the same size, the \textbf{addition/subtraction} is defined as $A \pm B = (a_{i,j} \pm b_{i,j})_{i,j=1}^{m \times n}$;
        \item for matrix $A \in \mathbb{F}^{m \times n}$ and scalar $\lambda \in \mathbb{F}$, the \textbf{scalar multiplication} is defined as $\lambda A = (\lambda a_{i,j})_{i,j=1}^{m \times n}$;
        \item for matrices\footnote{ Note that the number of columns in $A$ must be equal to the number of rows in $B$.} $A \in \mathbb{F}^{m \times n}$ and $B \in \mathbb{F}^{n \times p}$, the \textbf{matrix multiplication} $C := AB$ is a matrix $C \in \mathbb{F}^{m \times p}$ with elements $c_{i,j} := \sum_{\ell=1}^{n} a_{i,\ell}b_{\ell,j}$.
        \item for matrix $A \in \mathbb{F}^{m \times n}$, the \textbf{transpose} is a matrix $A^{\top} \in \mathbb{F}^{n \times m}$ with elements given by $A^{\top}_{i,j} = A_{j,i}$.
    \end{itemize}
\end{definition}

\begin{example}
    Suppose $\mathbb{F} = \mathbb{R}$. Then, consider 
    \begin{xequation*}
        A = \begin{bmatrix}
            1 & 1 & 2 \\
            2 & 2 & 1
        \end{bmatrix} \in \mathbb{R}^{2 \times 3}, \quad B = \begin{bmatrix}
            2 & 1 \\
            1 & 3 \\
            1 & 1
        \end{bmatrix} \in \mathbb{R}^{3 \times 2}
    \end{xequation*}

    We cannot add $A$ and $B$ since they have different sizes. However, we can multiply them:
    \begin{xequation*}
        AB = \begin{bmatrix}
            5 & 6 \\ 7 & 9
        \end{bmatrix}, \quad BA = \begin{bmatrix}
            4 & 4 & 5 \\
            7 & 7 & 5 \\
            3 & 3 & 3
        \end{bmatrix}
    \end{xequation*}

    To see why, for example, the upper left element of $AB$ is $5$, we can calculate it as $\sum_{\ell=1}^3 a_{1,\ell}b_{\ell,1} = 1 \times 2 + 1 \times 1 + 2 \times 1 = 5$.
\end{example}

\begin{exercise}
    Let $\mathbf{x} = \begin{bmatrix}
        1 \\ 2 \\ 1
    \end{bmatrix}$. Compute $\mathbf{x}^{\top}\mathbf{x}$ and $\mathbf{x}\mathbf{x}^{\top}$ using definitions above.
\end{exercise}

Notice that with the definition of matrix multiplication above, the matrix-vector product $A\mathbf{x}$ of $A \in \mathbb{F}^{m \times n}$ and $\mathbf{x} \in \mathbb{F}^n$ is
\begin{equation*}
    \begin{bmatrix}
        a_{1,1} & a_{1,2} & \dots & a_{1,n} \\
        a_{2,1} & a_{2,2} & \dots & a_{2,n} \\
        \vdots & \vdots & \ddots & \vdots \\
        a_{m,1} & a_{m,2} & \dots & a_{m,n}
    \end{bmatrix}\begin{bmatrix}
        x_1 \\ x_2 \\ \vdots \\ x_n
    \end{bmatrix} = \begin{bmatrix}
        a_{1,1}x_1 + a_{1,2}x_2 + \dots + a_{1,n}x_n \\
        a_{2,1}x_1 + a_{2,2}x_2 + \dots + a_{2,n}x_n \\
        \vdots \\
        a_{m,1}x_1 + a_{m,2}x_2 + \dots + a_{m,n}x_n
    \end{bmatrix},
\end{equation*}

which is (despite the scary look) a very expected expression.
For example, consider the following linear system of equations:
\begin{equation*}
    \begin{cases}
        x_1 + 2x_2 + 3x_3 = 2 \\
        2x_1 + x_2 + 3x_3 = 1 \\
        -x_1 + x_2 + 1x_3 = 0
    \end{cases}
\end{equation*}

By introducing\footnote{Despite writing $\mathbf{x}=(x_1,\dots,x_n)$, we consider vectors \emph{to be written in columns}, \textit{i.e.}, $\mathbf{x}=\begin{bmatrix}
    x_1 \\ \vdots \\ x_n
\end{bmatrix}$.}, $\mathbf{x}=(x_1,x_2,x_3)$, we can write it more succinctly as
\begin{equation*}
    A\mathbf{x} = \mathbf{t}, \quad A = \begin{bmatrix}
        1 & 2 & 3 \\
        2 & 1 & 3 \\
        -1 & 1 & 1
    \end{bmatrix}, \quad \mathbf{t} = \begin{bmatrix}
        2 \\ 1 \\ 0
    \end{bmatrix}.
\end{equation*}

\begin{remark}
    Among the provided operations above, the matrix-matrix multiplication seems the least intuitive.
    Here is the sketch why it is defined the way it is. 
    Recall that our goal is to study the operation $T: \mathbf{x} \mapsto A\mathbf{x}$.
    Now assume that we have another operation of such kind: $L: \mathbf{z} \mapsto B\mathbf{z}$.
    Naturally, we want to be able to compose these two operations; say, let
    \begin{equation*}
        F := T \circ L \qquad \textcolor{gray}{\textit{// where $(f \circ g)(x) = f(g(x))$.}}
    \end{equation*}
    
    Naturally, we expect that $F$ must also be of form $\mathbf{x} \mapsto C\mathbf{x}$.
    The only way for $C=AB$ to be correctly computed from $A$ and $B$ is by means of \Cref{def:matrix-ops}.
\end{remark}

Note that in \Cref{subsec:linalg:motivation} while considering simple systems of linear equations, we put a batch of (vertical) vectors into the matrix.
More concretely, given $\mathbf{v}_1,\dots,\mathbf{v}_n \in \mathbb{F}^n$, we formed the matrix $A := (\mathbf{v}_1,\dots,\mathbf{v}_n) \in \mathbb{F}^{n \times n}$.
Depending on what portion of $\mathbb{F}^n$ these vectors spanned, we drew different conclusions about the set of solutions of a linear equation.
Now we can naturally define the measure that will show this ``portion'' of $\mathbb{F}^n$.

\begin{definition}
    The \textbf{column rank} of a matrix $A \in \mathbb{F}^{m \times n}$ is the maximum number of linearly independent columns in $A$. 
    That is, if we represent the matrix in the form $A=(\mathbf{v}_1,\dots,\mathbf{v}_n) \in \mathbb{F}^{m \times n}$, then
    \begin{equation*}
        \text{col-rank}(A) := \dim \text{span}(\mathbf{v}_1,\dots,\mathbf{v}_n).
    \end{equation*}
\end{definition}

However, nothing restricts us from considering the matrix in the row-vector format.
That is, for $A \in \mathbb{F}^{m \times n}$, let
\begin{equation*}
    A = \begin{bmatrix}
        \mathbf{u}_1^{\top} \\
        \vdots \\
        \mathbf{u}_m^{\top}
    \end{bmatrix} \; \text{for some} \; \mathbf{u}_1,\dots,\mathbf{u}_n \in \mathbb{F}^n.
\end{equation*}

Similarly, we might ask: what portion of $\mathbb{F}^n$ do $\mathbf{u}_1,\dots,\mathbf{u}_m$ span?

\begin{definition}
    The \textbf{row rank} of a matrix $A \in \mathbb{F}^{m \times n}$ is the maximum number of linearly independent rows in $A$. 
    That is, if $A$ consists of rows $\mathbf{u}_1^{\top}, \dots, \mathbf{u}_m^{\top}$, then
    \begin{equation*}
        \text{row-rank}(A) := \dim \text{span}(\mathbf{u}_1,\dots,\mathbf{u}_m).
    \end{equation*}
\end{definition}

Fortunately, thanks to the following lemma, it turns out that it does not matter which type of rank we refer to.

\begin{lemma}
    The row rank and column rank of a matrix are always equal.
\end{lemma}

This allows us to simply define the \textbf{rank} of the matrix $A$ (and denote it by $\text{rank}(A)$) as either a row or column rank of $A$.

\subsection{Inner Product Spaces}

Now we switch to more geometric notions.
Commonly, when we need to work with the real field (for instance, when further dealing with lattices), we naturally want to give vectors a measure of length.
Additionally, we could ask the questions like ``what is the angle between these two vectors.''
Finally, with proper definitions, we can ask much more abstract questions: for example, given some continuous function $f: [0,1] \to \mathbb{R}$, can we define the ``length'' of this function?

To equip our vector space $V$ with a geometry, we turn it into the so-called \textbf{inner product space} which we define below.

\begin{definition}
    An \textbf{inner product} on $V$ is a function $\langle \cdot, \cdot \rangle: V \times V \to \mathbb{F}$ that satisfies the following properties:
    \begin{itemize}
        \item \textbf{positivity}: $\langle \mathbf{v}, \mathbf{v} \rangle \geq 0$ for all $\mathbf{v} \in V$;
        \item \textbf{definiteness}: $\langle \mathbf{v}, \mathbf{v} \rangle = 0$ iff $\mathbf{v}=0$;
        \item \textbf{additivity}: $\langle \mathbf{u}+\mathbf{v}, \mathbf{w} \rangle = \langle \mathbf{u}, \mathbf{w} \rangle + \langle \mathbf{v}, \mathbf{w} \rangle$ for all $\mathbf{u}, \mathbf{v}, \mathbf{w} \in V$;
        \item \textbf{homogeneity}: $\langle \lambda\mathbf{u}, \mathbf{v} \rangle = \lambda\langle \mathbf{u}, \mathbf{v} \rangle$ for all $\lambda \in \mathbb{F}$ and $\mathbf{u},\mathbf{v} \in V$.
        \item \textbf{symmetry}: $\langle \mathbf{u}, \mathbf{v} \rangle = \overline{\langle \mathbf{v}, \mathbf{u} \rangle}$ for all $\mathbf{v}, \mathbf{u} \in V$.\footnote{Here, $\overline{\lambda}$ denotes conjugation; for $\mathbb{F}=\mathbb{R}$, one can forget about it as $\overline{\lambda}=\lambda$.}
    \end{itemize}

    An \textbf{inner product space} is a vector space $V$ along with an inner product.
\end{definition}

\begin{remark}
    Note that the definition above does not make sense for finite fields or more generally for $\mathbb{Z}_N$ since expression $x \geq 0$ is meaningless for such structures.
    Consequently, when working with inner products over $\mathbb{Z}_N$, we typically require only \textit{additivity}, \textit{homogeneity}, and \textit{symmetry}.
\end{remark}

Plenty of functions can be built that satisfy the inner product definition, we will most commonly use the one that is usually called \textbf{dot product}.
\begin{example}
    Consider the vector space $V := \mathbb{F}^n$. The \textbf{dot product} on $\mathbb{F}^n$ is a function $\langle \cdot, \cdot \rangle: V \times V \rightarrow \mathbb{F}$, 
    defined for every $\mathbf{u}, \mathbf{v} \in \mathbb{F}^n$ as follows:
    \begin{xequation*}
        \langle \mathbf{u}, \mathbf{v} \rangle := \mathbf{u}^{\top}\mathbf{v} = \sum_{i=1}^{n} u_i v_i
    \end{xequation*}

    For example, let $\mathbb{F}=\mathbb{R}$ and let $\mathbf{u}=(1, 2, 3)$, $\mathbf{v} = (2, 4, 3)$. Then,
    \begin{xequation*}
        \langle \mathbf{u}, \mathbf{v}\rangle = \sum_{i=1}^{3}u_iv_i = 2 \cdot 1 + 2 \cdot 4 + 3 \cdot 3 = 2 + 8 + 9 = 19
    \end{xequation*}
\end{example}

\begin{example}
    We say that the square matrix $Q \in \mathbb{R}^{n \times n}$ is \textit{positive-definite} if $\mathbf{x}^{\top}Q\mathbf{x}>0$ for any non-zero $\mathbf{x} \in \mathbb{R}^n$.
    Then, the following function $\langle \cdot, \cdot \rangle_Q: \mathbb{R}^n \times \mathbb{R}^n \to \mathbb{R}$ is a valid inner product if $Q=Q^{\top}$:
    \begin{equation*}
        \langle \mathbf{x}, \mathbf{y} \rangle_Q := \mathbf{x}^{\top}Q\mathbf{y}.
    \end{equation*}
\end{example}

The notion of inner product naturally gives the measure of length, which is called a \textit{norm}.
Additionally, from the high-school geometry, one might remember that the dot product (with a proper rescaling) represents the angle between vectors in the Euclidean space.
In particular, when the dot product is $0$, two vectors are perpendicular/orthogonal. 
With our inner product definition, we can generalize this notion.

\begin{definition}\label{def:length}
    Let $V$ be an inner product space. Then,
    \begin{itemize}
        \item For $\mathbf{v} \in V$, the \textbf{norm} of $\mathbf{v}$, denoted $\|\mathbf{v}\|$, is defined by
        \begin{equation*}
            \|\mathbf{v}\| := \sqrt{\langle \mathbf{v}, \mathbf{v} \rangle}.
        \end{equation*}
        \item Two vectors $\mathbf{v},\mathbf{u} \in V$ are \textbf{orthogonal} if $\langle \mathbf{v}, \mathbf{u} \rangle = 0$.
    \end{itemize}
\end{definition}

\subsection{Other useful operations}

Finally, we list two other operations (\textit{Hadamard} and \textit{outer} products) that will be commonly used throughout.

\subsubsection{Hadamard Product}

Recall how inconvenient the matrix-matrix product looked like?
It turns out that a ``convenient'' way of doing this product (multiplying elementwise) has a special name.

\begin{definition}
    Suppose $A,B \in \mathbb{F}^{m \times n}$. The \textbf{Hadamard product} $A \odot B$ gives a matrix $C$ such that $C_{i,j} = A_{i,j}B_{i,j}$. Essentially, we multiply elements elementwise.
\end{definition}

\vspace{-4.2mm}

\begin{example}
    Consider $A = \begin{bmatrix}
        1 & 1 & 2 \\
        3 & 0 & 3
    \end{bmatrix}, B = \begin{bmatrix}
        3 & 2 & 1 \\ 0 & 2 & 1
    \end{bmatrix}$. Then, the Hadamard product is:
    \begin{xequation}
        A \odot B = \begin{bmatrix}
            1 \cdot 3 & 1 \cdot 2 & 2 \cdot 1 \\
            3 \cdot 0 & 0 \cdot 2 & 3 \cdot 1
        \end{bmatrix} = \begin{bmatrix}
            3 & 2 & 2 \\
            0 & 0 & 3
        \end{bmatrix}
    \end{xequation}
\end{example}

Despite the simplicity of such an operation, it is very common in various applications.
For example, suppose we want to prove that $\mathbf{z} \in \mathbb{F}^n$ is a bit-vector (\textit{i.e.}, every $z_i \in \{0,1\}$).
The natural way to represent this relation is as follows:
\begin{equation*}
    \mathbf{z} \odot (\mathbf{1}_n - \mathbf{z}) = \mathbf{0},
\end{equation*}

where $\mathbf{1}_n$ represents the vector of length $n$ consisting of $1$'s.
More generally, a R1CS relation, further considered in \Cref{section:r1cs}, is concisely represented in terms of Hadamard product between two specifically chosen large vectors.

\subsubsection{Outer Product}

The final product we introduce is the \textbf{outer product}.

\begin{definition}
    Given two vectors $\mathbf{u} \in \mathbb{F}^n$, $\mathbf{v} \in \mathbb{F}^m$ the \textbf{outer product} is a
    matrix whose entries are all products of an element in the first vector with an element 
    in the second vector:
    \begin{xequation}
        \mathbf{u} \otimes \mathbf{v} := \mathbf{u}\mathbf{v}^{\top} = \begin{bmatrix}
            u_1 v_1 & u_1 v_2 & \cdots & u_1 v_n \\
            u_2 v_1 & u_2 v_2 & \cdots & u_2 v_n \\
            \vdots & \vdots & \ddots & \vdots \\
            u_m v_1 & u_m v_2 & \cdots & u_m v_n
        \end{bmatrix}
    \end{xequation}
\end{definition}

\vspace{-4.2mm}

\begin{lemma}[Properties of outer product]
    For any scalar $c \in \mathbb{F}$ and $(\mathbf{u}, \mathbf{v}, \mathbf{w}) \in \mathbb{F}^n \times \mathbb{F}^m \times \mathbb{F}^p$:
    \begin{itemize}
        \item Transpose: $(\mathbf{u} \otimes \mathbf{v}) = (\mathbf{v} \otimes \mathbf{u})^{\top}$
        \item Distributivity: $\mathbf{u} \otimes (\mathbf{v} + \mathbf{w}) = \mathbf{u} \otimes \mathbf{v} + \mathbf{u} \otimes \mathbf{w}$
        \item Scalar Multiplication: $c(\mathbf{v} \otimes \mathbf{u}) = (c\mathbf{v}) \otimes \mathbf{u} = \mathbf{v} \otimes (c\mathbf{u})$
        \item Rank: the outer product $\mathbf{u} \otimes \mathbf{v}$ is a rank-1 matrix if $\mathbf{u}$ and $\mathbf{v}$ are non-zero
        vectors
    \end{itemize}
\end{lemma}

\vspace{-4.2mm}

\begin{example}
    Let $\mathbf{u}, \mathbf{v}$ be vectors over the real number $\mathbb{R}$, where
    $\mathbf{u}=(1, 2, 3)$ and $\mathbf{v} = (2, 4, 3)$. Then: 
    \begin{xequation*}
        \mathbf{u} \otimes \mathbf{v} = \mathbf{u}\mathbf{v}^{\top} = \begin{bmatrix}
        1 \\ 2 \\ 3
        \end{bmatrix}\begin{bmatrix}
        2 & 4 & 3
        \end{bmatrix} = \begin{bmatrix}
            1 \cdot 2 & 1 \cdot 4 & 1 \cdot 3 \\
            2 \cdot 2 & 2 \cdot 4 & 2 \cdot 3 \\
            3 \cdot 2 & 3 \cdot 4 & 3 \cdot 3
        \end{bmatrix} = \begin{bmatrix}
            2 & 4 & 3 \\
            4 & 8 & 6 \\
            6 & 12 & 9
        \end{bmatrix}
    \end{xequation*}
    
    Additionally, as we can see, the rows \#2 and \#3 in the resultant matrix can be represented
    as a linear combination of the first row, specifically by multiplying it by 2 and 3, 
    respectively. The same property applies to the columns. This demonstrates the property of the
    outer product, that the resulting matrix has a rank of 1.
\end{example}

\subsection{Exercises}

\textbf{Exercise 1.} Let $V := \{(x,y,z) \in \mathbb{R}^3: x+y+z=0\}$. Compute $\dim V$.
\begin{enumerate}[(A)]
    \item $0$.
    \item $1$.
    \item $2$.
    \item $3$.
\end{enumerate}

\textbf{Exercise 2.} Let $\mathbb{F}$ be a field and let $V := \{\mathbf{z} \in \mathbb{F}^n: \mathbf{z} \odot (\mathbf{1}_n - \mathbf{z}) = \mathbf{0}\}$.
When is $V$ a linear subspace of $\mathbb{F}^n$?
\begin{enumerate}[(A)]
    \item This never happens.
    \item $V$ is always a linear subspace of $\mathbb{F}^n$ with $\dim V = n$.
    \item $V$ is always a linear subspace of $\mathbb{F}^n$ with $\dim V = n-1$.
    \item Only when $\mathbb{F}=\mathbb{Z}_2$, otherwise it is not a linear subspace.
\end{enumerate}

\textbf{Exercise 3.} Let $A=\begin{bmatrix}
    1 & 2 \\ 0 & 1
\end{bmatrix}$. Compute $A^n$ for arbitrary $n \in \mathbb{N}$.
\begin{enumerate}[(A)]
    \item $\begin{bmatrix} 1 & 2n \\ 0 & 1 \end{bmatrix}$.
    \item $\begin{bmatrix} 1 & 2^n \\ 0 & 1 \end{bmatrix}$.
    \item $\begin{bmatrix} 1 & n \\ 0 & 1 \end{bmatrix}$.
    \item $\begin{bmatrix} n & 2n \\ 0 & n \end{bmatrix}$.
\end{enumerate}

\textbf{Exercise 4.} How many distinct $1$-dimensional linear subspaces (lines through the origin) exist in the vector space $\mathbb{F}_3^3$ over the finite field $\mathbb{F}_3$?
\begin{enumerate}[(A)]
    \item $9$.
    \item $13$.
    \item $26$.
    \item $27$.
\end{enumerate}

\textbf{Exercise 5*.} Recall that $V := \mathbb{R}[x]^{\leq 2}$ is a vector space over $\mathbb{R}$.
Define the inner product over $V$ as follows:
\begin{equation*}
    \langle p(x), q(x) \rangle := \int_0^1 p(x)q(x)dx.
\end{equation*}
Compute the ``length'' of the polynomial $x$ in $V$ (\textit{i.e.}, its norm $\|x\|$).
\begin{enumerate}[(A)]
    \item $1/\sqrt{3}$.
    \item $1/\sqrt{2}$.
    \item $\sqrt{2}$.
    \item $1$.
\end{enumerate}

\end{document}
